% !TEX root = ../main.tex

\chapter{真机实验}

\section{本章目标}

本章的重点是从部署流程、平台接口、运行模式、数据记录与安全机制等维度，对真机实验进行系统梳理。

对足式平台而言，真正的部署工作至少包含四类内容：首先，高层 NoMaD
推理必须在边缘计算设备上稳定加载和推理；其次，运动控制桥接必须能够把高层 waypoint 转化为底层可接受的速度命令；再次，系统应提供导航模式、探索模式以及多目标任务的
组织入口；最后，实验过程还应具备可视化、图像保存、日志记录和安全停机能力。

从更广泛的研究背景来看，仿真到真机的迁移（sim-to-real transfer）是具身智能领域公认的核心挑战之一。即便仿真环境中已经取得了理想的闭环表现，真实部署仍然
面临多重不确定性：视觉观测方面，真实光照变化、动态障碍物和纹理差异会导致视觉编码器的特征分布偏移；通讯方面，无线网络延迟抖动和丢包可能破坏控制闭环的时序假设；运动方面，真
实地面的摩擦系数变化、坡度和不平整性都会改变机器人的实际位移响应。对于本文所采用的基于扩散策略的视觉导航系统而言，最容易受到真实条件影响的环节有三个：一是视觉编码器对分布
外图像的鲁棒性，二是扩散采样延迟对控制周期的压力，三是高层 waypoint
到底层速度命令之间的映射精度。因此，本章的实验设计不仅要验证系统``能否运行''，还需要识别这些薄弱环节在真实条件下的实际表现。基于这一目标，本章内容先说明系统结构，再说明如何部署和验证，最后说明真实导航任务如何组织。

除上述技术层面的挑战外，足式机器人的真机部署还涉及若干仿真中不存在的工程约束。例如，真实机器人的电池续航决定了单次实验的最长运行时间；机器人关节的温度累积限制了连续高强度
运动的持续性；实验场地的空间限制和安全要求（如不能碰撞墙壁、不能跌落台阶）对导航策略的容错能力提出了更高要求。这些约束共同构成了本章实验设计的背景条件，也是后续结果分析时
需要考虑的因素。

\begin{figure}[htbp]
    \centering
    \includegraphics[width=0.8\linewidth]{shuanglianlu.png}
    \caption{Orin 与 Lite3 真机部署总体架构图}
    \label{fig:shuanglianlu}
\end{figure}

\section{真机系统的双链路设计}

为了把仿真阶段形成的视觉导航闭环迁移到 Lite3 真机，本文并不是简单地把 NoMaD 模型放到 Orin 上运行，而是将系统实现整理为两条相互交叉的链路。第一条是\textbf{横向实时推理链路}，描述数据如何在一个控制周期内从相机图像流动到速度指令；第二条是\textbf{纵向分层控制链路}，描述不同软件层和硬件层如何分工协作。前者实现一帧图像变成一条 twist 指令，后者实现这条指令安全、连续地进入 Lite3 运控。

\subsection{横向实时推理链路：从图像到 twist 指令}

横向链路是每个高层导航周期内实际发生的数据流，可概括为
\begin{equation}
\begin{aligned}
\text{Camera Images} &\rightarrow \text{Visual Encoder}
\rightarrow \text{Distance Predictor} \\
&\rightarrow \text{Diffusion Policy}
\rightarrow \text{PD Mapping}
\rightarrow \text{Twist Command}.
\end{aligned}
\end{equation}
其中，相机图像首先经过统一预处理，形成 NoMaD 所需的历史观察张量。当前默认配置中 \texttt{context\_size=3}，因此模型每次接收当前帧及其之前三帧，共四帧第一视角图像；导航模式还会输入一张实时子目标图像，并关闭 goal mask，使目标 token 真实参与条件编码。

视觉编码器是横向链路的第一段核心模型。它把四帧观察图像编码为观察 token，并把当前帧与目标图像拼接后编码为目标 token，再通过 Transformer 融合为统一条件特征。这个条件特征并不是只供扩散模型使用，而是同时喂给距离预测器和扩散策略头，因此视觉编码器的输出质量会同时影响局部定位和局部轨迹生成。

距离预测器位于视觉编码器之后，承担 topomap 局部定位功能。真机导航时，系统不会直接把最终目标图像作为唯一目标，而是在当前 \texttt{closest\_node} 附近构造滑动窗口，将窗口内每个候选节点分别与当前观察组成条件输入，预测当前状态到候选节点的离散距离。距离最小的候选节点用于更新定位指针，并进一步确定本轮扩散策略实际追踪的实时子目标。

扩散模型位于距离预测器之后，负责在实时子目标条件下生成未来局部 waypoint 序列。本文真机部署中采用 DDIM/DDPM 形式的轨迹去噪采样，输出的是机器人局部坐标系下的多步 waypoint，而不是底层关节动作。随后，中层 PD 映射根据选定 waypoint 计算前向速度、横向速度和偏航角速度，形成 Lite3 运动主机可接受的 twist 指令。横向链路由此完成从视觉输入到速度参考的端到端闭环。

\subsection{纵向分层控制链路：从 NoMaD 到运控}

纵向链路体现真机系统的工程分层，可概括为
\begin{equation}
\begin{aligned}
\text{High-level NoMaD Host}
&\rightarrow \text{Middle-level State Machine and PD Control} \\
&\rightarrow \text{Low-level UDP Bridge}
\rightarrow \text{Lite3 Locomotion Controller}.
\end{aligned}
\end{equation}
最高层是运行在 Jetson Orin 上的 NoMaD 导航主机，负责相机读取、topomap 管理、模型推理、任务队列、可视化、录像和日志。它只处理视觉导航意义上的目标、节点、waypoint 和任务状态，并不直接接触 Lite3 的关节级控制。

中层由状态机和 PD 控制映射组成。状态机负责组织 \texttt{idle}、\texttt{stand}、\texttt{navigate}、\texttt{explore}、\texttt{keyboard}、\texttt{estop} 等运行状态，决定何时进入导航、何时停止、何时触发急停或恢复；PD 映射则把扩散模型输出的局部 waypoint 转换为速度参考，并施加线速度、横向速度和角速度限幅，避免生成式策略偶发输出过大动作。

下层是真机桥接和 Lite3 运动主机。桥接层通过 UDP 与 Lite3 通讯，维护心跳和后台速度刷新线程，并将中层速度参考拆分为运动主机能够识别的底层速度通道。Lite3 自身的运控系统负责步态、平衡和关节级执行。这样，高层 NoMaD 可以保持与仿真阶段一致的输入输出接口，而真实机器人所需的通讯、起立、模式切换、心跳和安全停机都被封装在中下层。

两条链路在运行时并不是彼此独立的。横向链路在每个高层周期产生新的 twist 参考，纵向链路负责保证该参考在真实平台上被连续、安全地下发；当高层推理频率低于底层速度刷新频率时，桥接层会在相邻两次 NoMaD 推理之间持续发送最近一次限幅后的速度命令。正是这种横向端到端感知决策与纵向分层执行保护的结合，使得本文系统既能保留 NoMaD 的视觉导航能力，又能满足 Lite3 真机部署的实时性和安全性要求。

\section{Orin 真机部署架构}

\subsection{边缘计算平台的角色}

在本文部署方案中，Jetson Orin 承担高层导航主机的角色，即在其上完成相机读取、NoMaD
模型加载、扩散采样推理、状态机调度、导航主机交互和图像记录等工作。之所以选择 Orin 这类边缘计算平台，是因为它能够在机载条件下提供相对完整的 GPU
推理能力，从而使高层视觉导航不必完全依赖远程服务器。

从硬件规格来看，Jetson Orin 系列提供了多种配置等级，不同模块在 GPU 规模、统一内存容量和可配置功耗上并不完全相同。本文不把未在实验记录中逐项确认的核心数、Tensor Core 数量或持续功耗写作结论，而是把板端是否能够识别 \texttt{cuda} 设备、是否能稳定加载模型、以及 benchmark 中得到的端到端推理频率作为部署验收依据。相比之下，桌面级 RTX 4090 更适合离线训练和统计实验，但其供电、散热和体积条件并不适合直接随 Lite3 移动部署。

与其他可选的边缘计算平台相比，Jetson Orin 的优势主要体现在 CUDA 生态兼容性上。例如，Raspberry Pi 4/5 虽然功耗更低（约 5--
15\,W），但不具备 CUDA 推理能力，无法直接运行 PyTorch GPU 模型；Intel NUC 系列提供了较强的 CPU 算力，但缺乏高效的 GPU
推理后端，纯 CPU 推理下 NoMaD 的运行频率将远低于实时闭环要求；Google Coral 等 TPU 加速器虽然能效比高，但对 PyTorch
模型的支持需要额外的模型转换步骤，且对扩散模型的迭代采样支持不够成熟。综合考虑推理能力、功耗、生态兼容性和机载可行性，Jetson Orin
是当前足式机器人视觉导航部署中较为合理的选择。

但与此同时，Orin 也带来了明确约束。Orin 的算力与显存有限，因此高层推理、图像显示、图像保
存和控制桥接必须争夺同一套本地资源。这意味着真机部署必须从一开始就区分``调试模式''和``正式运行模式''：调试时允许开启可视
化和更多日志，正式运行时则优先保障闭环控制稳定性。

\subsection{真机执行链路与资源约束}

在双链路设计的基础上，Orin 真机部署进一步落实为可执行的软件链路：相机图像进入高层 NoMaD 推理模块，模型输出 waypoint 后交给中层 waypoint-to-command 映射，随后通过 Lite3 UDP 控制桥接进入运动主机。这一工程链路可以概括为

\begin{equation}
\text{Camera} \rightarrow \text{NoMaD Inference} \rightarrow \text{Waypoint-to-Cmd} \rightarrow \text{UDP Bridge} \rightarrow \text{Lite3}.
\end{equation}

需要进一步说明的是，真机相机并不要求原始输出分辨率与视觉编码器输入尺寸完全一致。当前部署链路会在推理前自动将图像缩放到策略配置中的
\texttt{image\_size}，例如 EfficientNet-B0、ConvNeXt-Tiny 与 ResNet-50 默认使用 $96\times96$
输入，DINOv2-Small 使用 $98\times98$ 输入。因此，对真机部署更关键的要求不是相机必须直接输出编码器尺寸，而是预处理过程与训练配置保持一致。在此基础上，相机标定应被视为可选增强步骤而非绝对前置条件：本文所使用的相机并未经过标定。


图像预处理流水线是连接相机原始输出与视觉编码器输入的关键环节。当前实现的预处理步骤包括：（1）从 USB 相机读取 $640\times480$ 的 BGR
图像；（2）将 BGR 格式转换为 RGB 格式；（3）将图像缩放至策略配置指定的输入尺寸（如 $96\times96$ 或
$98\times98$）；（4）将像素值归一化到 $[0,1]$ 范围并按照 ImageNet
均值和标准差进行标准化处理，使其与训练阶段的数据增强流水线保持一致。整个预处理过程在 CPU 上完成，耗时通常不超过 2\,ms，相对于推理耗时可忽略不计。

在推理效率优化方面，Orin 平台还具备从 PyTorch 向 TensorRT 转换的潜力。TensorRT 是 NVIDIA
提供的高性能推理引擎，能够通过算子融合、精度降低（FP16/INT8）和内存优化等手段显著提升推理速度。对于 NoMaD 的视觉编码器部分，理论上可以通过
\texttt{torch2trt} 或 ONNX 导出路径将其转换为 TensorRT 引擎，从而将视觉编码耗时从当前的约 57\,ms
进一步压缩。然而，扩散采样模块由于其迭代特性和动态控制流，TensorRT 转换的难度较高，目前本文暂未实施这一优化。未来若需要在更高控制频率下运行，TensorRT
优化将是一个值得探索的方向。

在真机部署阶段，高层 NoMaD 仍保持与仿真阶段一致的策略输入输出定义，而低层机器人通讯和步态细节则被吸收到桥接层与运动主机之中。从系统工程角度看，并不是另起一套真机专用算法，而是在仿真闭环系统的基础上切换到真实后端。

\subsection{真机桥接层}

Lite3 的运动主机并不直接接受扩散轨迹或视觉目标图像这类高层抽象概念，它接受的是速度指令、心跳信息和模式切换命令。因此，在 Orin
和机器人之间必须存在一个真机桥接层，用来完成三项关键任务：其一，把高层 motion command 映射为底层速度控制命令；其二，维护与运动主机之间的 UDP
通讯与心跳；其三，对外暴露统一平台接口，使导航主机与状态机能够在不理解底层协议细节的前提下驱动真机。

在 UDP 通讯的具体实现中，本文并不是把三维速度、心跳计数器和时间戳打包成一个自定义综合数据包，而是复用 Lite3 上位机控制接口提供的命令封装。桥接层中的 \texttt{Lite3Controller} 会启动接收线程和自动心跳线程；当前配置下心跳线程以 \texttt{5\,Hz} 发送心跳，速度控制线程以 \texttt{25\,Hz} 重复下发最近一次高层速度参考。由于 Lite3 协议中的速度轴命令是按前后、左右和偏航三个通道分别发送的，\texttt{send\_twist()} 实际会把一个高层三维速度向量拆成三条底层速度命令。高层 NoMaD 推理频率低于这一速度刷新频率，因此桥接层会在两次高层推理之间持续发送最近一次限幅后的速度参考，以满足运动主机对速度命令连续性的要求。

从端到端延迟预算的角度分析，完整的控制链路延迟可以分解为以下几个环节：相机图像采集约 25--40\,ms（取决于曝光时间和 USB 传输），图像预处理约
2\,ms，视觉编码约 57\,ms，扩散采样约 87\,ms，waypoint-to-cmd 映射不超过 1\,ms，UDP 发送与网络传输约 1--
3\,ms。因此，从相机捕获图像到机器人收到对应速度命令的总延迟约为 170--190\,ms，对应约 5--6\,Hz 的闭环控制频率。对于当前保守低速行走场景（约 0.12\,m/s），这一延迟水平是可接受的；但若未来需要支持更高速度的运动（如 0.8\,m/s 以上），则需要通过减少 DDIM 步数或引入 TensorRT
优化来进一步压缩推理延迟。

通讯故障处理是桥接层设计中不可忽视的环节。当前代码中的保护更偏向``上层主动检测与软急停''，而不是在桥接层内部维护一个 500\,ms 发送失败计时器。具体来说，UDP 发送异常会被记录并限频告警；交互式导航主机在真机后端会周期性检查 \texttt{is\_fallen()} 和 \texttt{is\_connection\_stale(timeout=2.0)}，当机器人摔倒或运动主机状态超过 2\,s 未更新时，自动调用 \texttt{emergency\_stop()} 并退出当前任务。任务正常结束时，系统通过 \texttt{controlled\_stop()} 发送零速度并释放手动控制；任务失败或急停时，则调用桥接层的 \texttt{soft\_estop()}。

这层桥接设计与第四章中的平台抽象接口是一致的。它既是``同一高层架构、双后端复用''在工程上的落地形式，也是真机部署安全性的关键保障。没有桥接层，导航主机将不得不直接操作底层通讯协议，这会使系统既难以维护，也难以在不同硬件平台之间扩展。

\subsection{真机部署中的时序等待机制}

真机部署中的时序等待并不是算法层面的``暂停思考''，而是围绕相机感知、运动主机、UDP 通讯和人机交互所建立的工程时序保护。与 MuJoCo 仿真不同，真实平台中的相机打开、机器人起立、模式切换和速度指令下发均需消耗物理时间；若上层程序在这些环节不等待或等待不足，则可能出现相机帧尚未稳定、机器人尚未站稳、运动主机尚未完成模式切换而高层策略已经开始下发速度的情况。因此，本文将这些等待机制视为真机系统可靠性的一部分，而非临时调试代码。

为便于讨论，本节先在系统层面抽象出四类接口：相机感知接口、平台运动接口、通讯桥接接口与人机交互接口。相机感知接口对外提供稳定的图像帧序列；平台运动接口对外提供起立、模式切换与停机等动作请求；通讯桥接接口将上层速度参考转化为底层连续指令流；人机交互接口处理操作者输入并维持空闲状态下的命令轮询。每一类接口都对应一组特定的时序约束，下面分别说明。

第一类等待存在于相机感知接口中。当 USB 或 CSI 相机首次打开失败时，系统以约半秒的间隔进行最多三次重试；首次成功打开后，还会丢弃数帧作为预热。该设计的目的是避开 USB 设备刚释放、刚插入或曝光自动调整尚未稳定时的异常帧，确保进入推理模块的图像分布与训练数据保持一致。

第二类等待存在于平台运动接口的起立与模式切换流程中。当向运动主机发送起立指令后，机器人需要约三秒完成站立姿态稳定；随后向自主控制模式与移动控制模式的切换分别需要约半秒的固定等待。上层接口允许操作者通过单一参数指定起立的总等待时间，若该值大于底层默认等待，则在底层等待之外补足相应时长。本文实验中将起立总等待设置为三秒，即不在底层之外额外延长；当现场观察到机器人起立后仍有明显晃动时，可将该参数适当增大以延长稳定时间。

第三类等待用于维持实时控制节拍。Lite3 通讯协议要求速度轴指令持续下发，因此通讯桥接接口在后台维持两条独立线程：速度刷新线程以约 25\,Hz 的频率重复下发最近一次速度参考，对应约 40\,ms 的循环间隔；心跳线程以 5\,Hz 的频率发送心跳信号，对应约 200\,ms 的循环间隔。这两个频率均满足运动主机对底层速度命令和心跳频率的最低要求。由于高层 NoMaD 推理频率约为 6--7\,Hz，明显低于底层速度刷新频率，桥接接口会在相邻两次高层 waypoint 更新之间持续重复最近一次限幅后的速度命令，使运动主机看到连续控制流而非稀疏突发指令。

第四类等待存在于人机交互接口的空闲与遥操作状态中。当系统处于空闲状态等待新命令时，导航主机以约 50\,ms 的周期轮询终端输入并保持站立姿态；键盘遥操作模式采用同样的循环间隔，对应约 20\,Hz 的手动控制节拍。该等待量在保证主机能够及时响应起立、导航、探索与急停等在线命令的同时，避免空闲状态下 CPU 资源被无意义循环占满。对于真机安全而言，空闲循环中的短等待还与连接陈旧检测协同工作：当运动主机状态超过两秒未更新时，导航主机会触发自动退出或急停处理。


\section{真机部署流程}

\subsection{阶段一：Orin 端独立调试}

真机部署的第一阶段为 Orin 端独立调试，即在 Lite3 平台未接入的条件下，先验证相机读取、模型加载、CUDA 环境、端到端推理频率以及图像保存链路是否满足部署要求。该阶段的目的在于将高层推理子系统与机器人控制子系统的故障域加以隔离：当高层推理本身存在异常时，过早接入真机会使两个子系统的问题相互掩盖，进而显著增加故障定位的复杂度。

阶段一中需要依次验证的内容包括：相机能否稳定提供图像帧；NoMaD 模型权重能否在 Orin 上完成加载并被正确分配到 GPU；统一推理模块能否在板端输出合理的 waypoint；DDIM 等参数下的端到端推理频率是否处于可用区间；以及结果保存目录能否正常生成。从研究意义上看，该阶段不仅验证了软件链路的连通性，更确认了第四章建立的统一推理模块在边缘计算平台上的可迁移性。

相机标定的处理方式同样在阶段一中确定。当相机原始画面中的直线结构基本正常、边缘畸变较小时，可直接进入独立调试，因为部署链路会自动完成输入尺寸对齐；当相机存在明显桶形畸变或针孔模型偏差时，则需先通过离线标定获得相机内参与畸变参数，再在 Orin 端启用桥接层去畸变。考虑到部署可用性与工程可维护性的平衡，本文将相机标定实现为可选启用的桥接层功能，而非所有实验场景的强制前置步骤。

除畸变问题外，真机相机的宽高比同样需要单独说明。本文当前 USB 相机的原始输出分辨率为 $640\times480$，而 NoMaD 基线视觉编码器的输入尺寸为 $96\times96$。直接将 $4:3$ 图像缩放为 $1:1$ 会引入一定横向几何压缩，并非物理成像意义上的最真实变换；但从学习系统的角度看，只要训练阶段、topomap 节点、目标图像和实时观测均采用同一预处理方式，则该变换在输入分布一致性上仍然成立。当前系统将图像对齐方式作为推理模块参数，默认采用 \texttt{stretch} 以保留完整视场并与原始 NoMaD 训练预处理保持一致；同时提供 \texttt{center\_crop} 和 \texttt{letterbox} 两种可选方式，分别用于保持几何比例但裁剪视场，以及保留完整视场但引入黑边。本文在首轮真机部署中采用 \texttt{stretch}；若后续实验观察到横向距离估计或转弯幅度出现系统偏差，则可通过阶段一的对比测试对三种对齐方式的影响进行进一步评估。

\subsection{阶段一实测结果与分析}


\begin{table}[htbp]
    \centering
    \caption{Orin 阶段一独立调试实测结果}
    \label{tab:orin-stage1-results}
    \begin{tabular}{p{0.20\linewidth}p{0.24\linewidth}p{0.42\linewidth}}
        \toprule
        测试项 & 实测结果 & 结果说明 \\
        \midrule
        相机读取 & $640\times480$，约 $40.0\,\mathrm{ms}$，约 $25.0\,\mathrm{FPS}$ & 说明 USB
        相机链路稳定，能够持续提供导航输入图像。 \\
        模型加载 & 2.87\,s，设备为 \texttt{cuda} & 说明 Jetson Orin 集成 GPU 已被正确识别，NoMaD
        权重可在板端完成加载。 \\
        基准推理（DDIM-5） & 视觉编码 56.9\,ms；扩散采样 86.7\,ms；端到端 143.6\,ms；约 7.0\,Hz &
        说明高层推理模块在 Orin 上已达到可用于后续桥接验证的运行频率。 \\
        完整流水线 & 首轮 2895.5\,ms；后续约 151--157\,ms & 说明系统存在明显冷启动开销，但稳态运行阶段可保持约 6.4--
        6.6\,Hz。 \\
        Bridge 相机测试 & 连续 30 帧均成功读取 & 说明桥接层相机包装与独立相机测试结果一致。 \\
        Deployment Checklist & 关键文件均存在 & 说明阶段二联调所需的模型、主机、控制器和资产文件链路完整。 \\
        \bottomrule
    \end{tabular}
\end{table}

从结果上看，阶段一最重要的结论有三点。第一，相机链路与模型链路已经分别通过验证。第二，\texttt{benchmark} 测试给出的约 7.0\,Hz
更适合作为 Orin 高层推理能力的代表值，因为它排除了相机初始化和首轮图优化带来的额外开销；相比之下，\texttt{pipeline} 测试的原始平均值
427.3\,ms 虽然真实，但被第一轮 2895.5\,ms 的冷启动延迟显著拉高，不能直接代表系统进入稳态后的闭环能力。第三，完整流水线在首轮之后基本稳定在 151
--157\,ms 区间，说明当前 \texttt{DDIM-5} 配置下系统已经具备进入 Lite3 低速联调的实时性基础。

相机读取的 25\,FPS 意味着图像采集不会成为推理链路的瓶颈，因为高层推理频率仅约
7\,Hz，相机能够以远高于推理需求的速率供给图像帧。模型加载耗时 2.87\,s 属于一次性开销，在系统启动阶段完成，不影响运行时性能；但这一数值也表明，若系统需要在运
行中动态切换模型权重（例如从导航模型切换到探索模型），则需要额外考虑模型切换的延迟成本。基准推理中视觉编码占 56.9\,ms、扩散采样占 86.7\,ms
的分解数据表明，当前推理延迟的主要贡献来自扩散采样而非视觉编码，这与第三章中的消融实验结论一致，也进一步支持了``减少 DDIM 步数是降低推理延迟最有效手段''的判断。

因此，阶段一实验结果通过。这并不意味着真机闭环导航已经完成，而是意味着 Orin 端已经满足以下前置条件：一是图像输入稳定，二是模型能以
\texttt{cuda} 设备运行，三是高层输出 waypoint 数值合理且非零，四是部署资产链路完整。

\subsection{阶段二：Orin 与 Lite3 联合调试}

阶段二是 Orin 与 Lite3 联合调试阶段，包括网络连通验证、UDP 桥接测试、起立测试、低速速度测试、交互式导航测试和目标导航测试。该阶段采用从低风险动作向完整闭环渐进的组织原则：足式机器人各底层环节的不稳定性会在高层导航阶段被放大，将测试粒度按风险等级排序有助于在最早的子环节中暴露故障。

在具体顺序上，首先确认 Orin 与 Lite3 运动主机之间的网络与心跳是否正常；其次执行起立测试，验证机器人能否从趴卧状态平稳过渡到自主控制模式；随后进行低速 twist 测试，验证桥接层能否在固定低速下连续下发速度命令；最后再启动导航主机驱动的探索或目标导航流程。

阶段二各子步骤的具体内容如下。当前无线联调采用 Lite3 WiFi 热点，Orin 由热点自动分配到的地址为 \texttt{192.168.2.17}，用于操作者电脑通过 SSH 进入导航主机；Lite3 运动主机地址使用 \texttt{192.168.2.1}，用于桥接层发送 UDP 控制命令。二者的含义不能混淆：前者是导航主机地址，后者才是 \texttt{robot\_ip}。网络连通测试要求 Orin 与 Lite3 运动主机在同一局域网内能够双向 ping 通，且 UDP 端口未被防火墙阻断。心跳测试要求桥接层启动后，运动主机能够在日志中确认已收到有效心跳包。起立测试通过桥接层发送起立指令，观察机器人是否能从趴卧状态平稳过渡到站立姿态，该步骤同时验证了控制权限是否已正确移交给 Orin 端。低速 twist 测试以固定的小幅线速度（如 $v_x = 0.1$\,m/s）和零角速度驱动机器人直线前进 2--3 秒，用于验证速度命令的方向性和幅度合理性。交互式导航测试允许操作者通过键盘或命令行指定临时目标图像，观察机器人是否能朝目标方向产生合理运动响应。

阶段二实验结果通过，意味着 orin 与 Lite3 联合系统已经具备真机部署的前提条件。

\section{真实 topomap 采集与滑动窗口导航流程}

本节把本文已经实现的真机导航系统作为一个完整流程进行整理。与单张目标图像导航不同，当前系统在真实场景中采用\textbf{先采集有序 topomap，再沿 topomap 滑动推进}的方式完成目标导航。其核心思想是：先让操作者沿真实路线拍摄一串第一视角图像，形成从起点到终点的视觉路标序列；正式运行时，Lite3 从序列起点附近出发，系统在局部滑动窗口中用距离预测器估计当前最接近的节点，再把由该节点推进得到的实时子目标图像送入扩散策略，生成下一段局部 waypoint。机器人不断重复这一过程，直到定位指针推进到指定目标节点。

\subsection{导航输入与实时子目标}

真机端的高层算法仍然由 NoMaD 模型承担，但真实导航任务的输入组织与单张目标图像导航不同。系统并不是把终点图像一次性作为固定 goal 输入，而是根据真实 topomap 和滑动窗口机制，在每个高层周期选择一个实时子目标图像。

相机与预处理模块负责从 Orin 连接的 USB/CSI 相机中读取实时图像，并将图像转换为与训练配置一致的输入张量。本文当前默认配置中，\texttt{context\_size=3}，因此每次高层推理使用当前帧及其之前三帧，共四帧历史观察图像；若刚启动时缓存不足，系统会重复最早帧补齐上下文。导航模式下还会额外输入一张目标图像，但这张图并不是整段任务的最终终点图，而是本轮滑动窗口机制选出的实时子目标节点图像。此时 goal mask 关闭，目标 token 参与视觉条件编码；探索模式则使用被 mask 的目标条件。

视觉编码器负责把上述四帧观察图像和实时子目标图像编码为统一条件特征。在 NoMaD 的 ViNT 结构中，四帧观察图像分别经过观察编码器得到观察 token，当前帧与目标图像沿通道维度拼接后进入目标编码器得到目标 token，随后这些 token 经过 Transformer 自注意力融合，并池化为固定维度的条件向量。这个条件向量被后续两个头共享，因此视觉编码器、距离预测器和扩散策略头不是彼此独立的三段程序，而是共享同一个视觉条件空间。

距离预测器用于 topomap 局部定位。给定当前观察和滑动窗口中的若干候选节点图像，系统会把同一组观察帧分别与窗口内每一张候选图像组成输入，预测当前状态到候选节点的离散距离。距离最小的候选节点被视为当前视觉上最接近的位置，再经过反跳变限制更新为新的 \texttt{closest\_node}。这一距离不是物理米制距离，而是训练数据中相对导航步数意义下的距离估计。

扩散策略头负责生成局部 waypoint。系统不会把整个 topomap 或最终终点一次性送入扩散模型，而是只把滑动窗口选出的一个实时子目标节点作为目标条件。扩散策略在该条件特征上通过 DDIM/DDPM 去噪生成未来若干步局部 waypoint，随后中层控制器从预测轨迹中选取指定索引的 waypoint，将其转换为 Lite3 可执行的前向速度、横向速度和偏航角速度。最后，真机桥接层通过 UDP 将速度命令持续发送给 Lite3 运动主机，并在两次高层推理之间保持最近一次速度参考。

\subsection{真实 topomap 的采集与组织}

在正式导航之前，必须先在真实环境中采集 topomap。采集脚本会调用 Orin 相机沿路线保存图像。

采集方式分为自动和手动两种。自动模式按照设置的时间间隔保存图像，适合操作者能够以稳定速度沿路线移动相机的场景；手动模式程序会显示实时预览窗口，操作者每到一个希望记录的节点位置后保存当前节点。

采集得到的图像按照文件名排序后构成有序序列
\begin{equation}
\mathcal{M}=\{I_0,I_1,\ldots,I_N\},
\end{equation}
其中 $I_0$ 对应导航起点附近的第一视角，$I_N$ 通常对应默认终点。节点编号保证与真实路线的前进顺序一致，这是系统能够滑动推进的前提。

在节点间距上，过密的 topomap 会增加存储和匹配开销，但相邻节点高度相似时不一定显著提升效果；过疏的 topomap 则会造成相邻目标之间的视觉跳变，使模型难以在中间位置生成稳定轨迹。如室内走廊topomap，相邻节点间隔距离在1m至1.5m之间； 较空旷地面topomap，相邻节点间隔距离在0.5m至1m之间。

\subsection{滑动窗口导航机制}

真实导航开始时，任务队列将定位指针初始化为
\begin{equation}
c_0=0,
\end{equation}
即默认机器人从 topomap 的第一张图像附近出发。目标节点 $g$ 默认为序列最后一个节点 $N$，也可以通过指定 topomap 中某个图片文件名来提前结束任务。需要强调的是，$g$ 只是当前任务的终点上界，并不表示每个推理周期都直接把 $I_g$ 作为扩散目标。

在第 $t$ 个高层周期，系统以当前定位指针 $c_t$ 为中心构造局部候选窗口。当前代码对应的窗口边界为
\begin{equation}
s_t=\max(c_t-r,0),\qquad e_t=\min(c_t+r+1,g),
\end{equation}
候选集合实际为
\begin{equation}
\mathcal{W}_t=\{I_{s_t},I_{s_t+1},\ldots,I_{e_t}\}.
\end{equation}
其中 $r$ 为窗口半径参数。实际右边界使用 $c_t+r+1$，是为了在当前节点前后搜索的同时保留一个更靠前的候选节点，便于机器人接近当前节点后继续向下一站推进。起点附近窗口左端被裁剪到 0，终点附近窗口右端被裁剪到 $g$，因此窗口会先从起点附近展开，中途随 $c_t$ 向后滑动，最后停在目标节点 $g$ 之前或包含 $g$ 的局部范围内。

随后，系统把同一组观察帧分别与窗口内每个候选节点图像组成目标条件，调用视觉编码器和距离预测器得到一组距离估计：
\begin{equation}
d_i = D_\theta(O_t, I_i),\qquad i\in[s_t,e_t].
\end{equation}
距离最小的节点
\begin{equation}
\hat{c}_t=\arg\min_{i\in[s_t,e_t]}d_i
\end{equation}
表示当前图像在局部窗口内最接近的 topomap 节点。为了防止视觉误匹配导致定位指针突然跳远，实际实现时还会把该结果限制在 $[c_t-1,c_t+3]$ 范围内，并且不允许超过目标节点 $g$。这一限制使指针可以缓慢后退一格以修正误差，但最多只允许每轮向前推进三个节点。

扩散模型实际使用的目标节点记为 $u_t$。它并不总是窗口中距离最小的节点：如果距离预测器认为机器人已经足够接近当前最近节点，即 $d_{\min}<\tau$，系统会把扩散条件推进到该节点之后的下一张图像；否则仍使用当前最近节点。形式上可以概括为
\begin{equation}
u_t=
\begin{cases}
\min(\hat{c}_t+1,g), & d_{\min}<\tau,\\
\hat{c}_t, & d_{\min}\ge \tau.
\end{cases}
\end{equation}
因此，距离预测器的作用是定位和推进指针，扩散策略的作用是朝实时子目标 $I_{u_t}$ 生成下一段局部轨迹。机器人执行该局部轨迹后，下一轮再读取新的相机帧、更新四帧上下文、重新构造窗口并重复上述过程。整个过程不是一次性规划全局轨迹，而是“观察--局部定位--选择子目标--生成 waypoint--执行--再观察”的闭环。

\section{导航模式、探索模式与图像记录机制}

\subsection{导航模式与探索模式}

真机端同样支持导航模式与探索模式两类高层运行方式。导航模式下，系统明确接收目标图像或目标节点，并由 NoMaD
生成朝向该目标的局部轨迹；探索模式下，系统弱化明确目标条件，使高层策略依赖当前视觉上下文和训练中学到的先验运动趋势继续前进。由于 goal mask
已经在训练阶段统一了这两种能力\cite{nomad2023}，因此真机端不需要切换完全不同的模型，只需要切换任务模式与目标输入即可。


\subsection{可视化模式与无可视化模式}

真机部署中另一个需要明确区分的问题是可视化模式与无可视化模式。可视化模式主要用于调试、录屏与实验观察，此时系统会显示实时相机画面、当前目标图像以及必要的状态信息；无可视化
模式则面向正式运行，优先保证闭环稳定与计算资源利用，不打开图形窗口，只保留日志与必要图像保存。

在Orin 这类边缘计算平台上，GUI 显示和图像刷新会消耗额外资源，从而影响推理周期稳定性。将``是否显示
图像''设计为可切换项，意味着系统可以在调试阶段最大化可观测性，而在正式运行阶段最大化稳定性。这也是本文在系统层特别强调``推理模块''和``部署模块''分离的原因之一。


\section{安全机制}

本文在部署链路中设置了多处可验证的安全入口，包括任务正常结束时的零速度停止、任务失败或急停时的软急停、导航主机对摔倒和连接陈旧状态的检测，以及实验者保留的遥控器或硬件断电手段。与此同时，桥接层持续维护底层心跳和后台速度刷新，减少高层推理周期抖动对运动主机的影响。

从层级化安全架构的角度，当前系统可以概括为四类保护。第一类是任务级停止：状态机在任务完成时调用 \texttt{controlled\_stop()}，真机后端会发送零速度并释放手动控制。第二类是软件急停：当状态机进入 \texttt{estop} 或任务失败路径时，会调用 \texttt{emergency\_stop()}，在 Lite3 真机后端对应 \texttt{soft\_estop()}。第三类是状态监测：导航主机在交互式运行中检查 \texttt{is\_fallen()}，同时以 2\,s 阈值检查运动主机状态是否陈旧，触发后会急停并退出当前任务。第四类是实验规程层面的人工保护：监护者应保留遥控器、物理急停或断电手段，用于覆盖软件链路之外的极端情况。

这些安全机制并不是与 NoMaD 无关的额外功能，而是高层生成式策略能够进入真机的先决条件\cite{diffusionpolicy2023,hwangbo2019anymal,lee2020learning}。因为生成式模型输出本身具有分布性，如果系统没有明确的速度限幅、恢复和急停通道，即使轨迹在大多数时刻合理，也可能在局
部异常时对平台造成风险。


\section{真机闭环实验结果}

本节给出 Lite3 真机平台上已经完成的三类实验结果：真实场景 topomap 采集、多场景目标导航与多场景探索。所有实验均在保守速度限幅下进行（前向线速度上限 $0.12\,\text{m/s}$、偏航角速度上限 $0.35\,\text{rad/s}$），并在每次运行前后保存 session 日志、运行视频和关键帧图像，以保证结果的可复现性。

\subsection{真实场景 topomap 采集}

正式导航实验前，本文按前述滑动窗口导航流程依次在多个真实场景中采集了有序 topomap 节点序列。采集完成后，每条 topomap 均通过人工回放检查序列的空间连续性与朝向一致性，确保进入正式导航实验时不会因节点排序异常导致定位指针推进错误。

\begin{figure}[htbp]
    \centering
    \includegraphics[width=0.8\linewidth]{map.png}
    \caption{真实场景 topomap 采集}
    \label{fig:map}
\end{figure}

\subsection{多场景目标导航实验}

为评估第三章离线实验筛选出的推理配置在真机上的实际表现，本文在前述真实 topomap 上组织了多场景目标导航实验，并对四组推理配置进行了闭环对比。四组配置分别为：基线 DDPM 配置（沿用 NoMaD 默认采样器与步数）、最快方案（DDIM-2、CFG=0、TTS-8）、最优方案（DDIM-2、CFG=2、TTS-8）和对比方案（DDIM-3、CFG=0、无 TTS）。每组配置在多场景中各运行若干次，并以``机器人沿 topomap 推进至目标节点且未触发急停或人工干预''作为成功条件。

实验结果如表~\ref{tab:real-multiscene-navigate} 所示。在当前真实场景规模与速度限幅下，四组配置在多场景目标导航中均达到 100\% 成功率。在此基础上，更具区分意义的指标是闭环频率，反映各推理配置在真实部署条件下的实时性差异。

\begin{table}[htbp]
    \centering
    \caption{多场景目标导航的真机闭环实验结果}
    \label{tab:real-multiscene-navigate}
    \begin{tabular}{lcc}
        \toprule
        推理配置 & 成功率 & 闭环频率/Hz \\
        \midrule
        基线 DDPM & 100\% & 2.89 \\
        最快方案（DDIM-2, CFG=0, TTS-8） & 100\% & 6.80 \\
        最优方案（DDIM-2, CFG=2, TTS-8） & 100\% & 3.43 \\
        对比方案（DDIM-3, CFG=0, 无 TTS） & 100\% & 5.90 \\
        \bottomrule
    \end{tabular}
\end{table}

从成功率维度看，四组配置在多场景目标导航中均达到 100\%。这一结果表明，在当前真实场景规模与保守速度限幅下，系统具有充分的鲁棒性裕度：即便最慢的基线 DDPM 配置（2.89\,Hz、约 346\,ms 单帧周期）也能稳定完成沿 topomap 推进的导航，CFG 与 TTS 等推理质量增强项的引入也未导致额外失败。因此，在当前实验范围内，成功率本身不再具有区分四组配置优劣的能力，更具区分意义的指标转向了闭环频率。

闭环频率与推理配置之间呈现出与第三章离线分析高度一致的趋势。基线 DDPM 使用完整 10 步去噪，单次扩散采样耗时最长，闭环频率最低（2.89\,Hz）；将采样器替换为 DDIM 后，去噪步数大幅减少：DDIM-3 的对比方案提升至 5.90\,Hz，DDIM-2 的最快方案进一步提升至 6.80\,Hz。这一``DDIM 步数减少 $\Rightarrow$ 闭环频率单调提升''的关系直接验证了第三章关于``少步 DDIM 是降低推理延迟最有效手段''的结论。CFG 的开销同样可在频率数据中清晰读出：最优方案与最快方案除 CFG 权重外其他配置完全一致，但最优方案（CFG=2）的 3.43\,Hz 约为最快方案（CFG=0）的 6.80\,Hz 的一半，原因在于 CFG 每个去噪步需执行条件与无条件两次前向传播，等效将扩散采样计算量翻倍，这与第三章和第四章中``CFG 在提升条件牵引力的同时引入约一倍的采样开销''的结论相吻合。TTS 的开销则相对可控：最快方案（DDIM-2+TTS-8，6.80\,Hz）的频率仍高于对比方案（DDIM-3 无 TTS，5.90\,Hz），表明 TTS 候选轨迹通过批量并行采样后，其相对单条轨迹的额外耗时主要落在轻量评分函数上，远低于扩散去噪本身，这与第三章关于``批量 TTS 在 GPU 上具有较好并行效率''的结论保持一致。综合来看，真机闭环频率与第三章离线消融实验给出的相对快慢顺序完全吻合，说明算法层结论在迁移到真实部署条件后并未发生倒置。

值得进一步分析的是，尽管基线 DDPM 仅以 2.89\,Hz 运行，机器人仍然能够稳定完成所有真机导航任务。其一，本文采用的保守速度限幅（前向 $0.12\,\text{m/s}$）使得机器人在两次相邻高层推理之间至多移动约 $0.12/2.89\approx 4.2\,\text{cm}$，远小于 topomap 的节点间距，因此低频高层更新不会在子目标层面引起明显的定位指针漂移。其二，通讯桥接层以 25\,Hz 的频率持续重复最近一次速度命令，使运动主机在两次 NoMaD 推理之间仍能看到连续控制流，避免了高层频率不足导致的速度指令稀疏问题。其三，滑动窗口 topomap 机制为高层导航提供了空间锚点，即使单次扩散采样存在轻微抖动，距离预测器在下一周期仍可将定位指针拉回正确节点。三者共同保证了高层频率的下限被推至 $\sim 3\,\text{Hz}$ 量级仍不破坏闭环可行性。

最后，最快方案在保持 100\% 成功率的同时将闭环频率从基线的 2.89\,Hz 提升至 6.80\,Hz，约为基线的 2.4 倍。这一结果具有双重意义：一方面，基线 DDPM 的成功部署证明了系统对低频高层推理具有充分容错能力；另一方面，最快方案带来的频率余量为后续支持更高速度运动和更复杂动态环境提供了方案。

需要说明的是，上述``100\% 成功率''是当前室内实验场景规模下的结论。受训练集、电池续航、场地尺寸与人工监护成本的限制，首先本文算法训练只用了室内数据集，其次本文真机实验的样本数有限；若将场景扩大到室外场景、含强光变化、动态障碍或更长走廊的环境中，不同推理配置之间的差异有望更显著，高频配置在面对快速变化场景时的优势也可能进一步放大。

\begin{figure}[htbp]
    \centering
    \includegraphics[width=0.5\linewidth]{real-robot-eg.png}
    \caption{真机目标导航过程}
    \label{fig:real-multiscene-navigate-placeholder}
\end{figure}

\subsection{多场景探索实验}

探索模式实验同样在多个真实场景中开展。与目标导航不同，探索任务不提供目标 token，而是依靠训练阶段习得的运动先验在当前视觉上下文中持续生成可行轨迹。实验过程中，机器人能够沿走廊或开阔区域稳定前进，在接近墙壁或障碍物前完成减速，并由扩散策略生成绕行或转弯的局部 waypoint。这一行为模式与第三章关于 goal mask 的分析一致，说明探索模式在真机平台上同样得到了行为层面的验证。

\section{本章小结}

本章围绕 Orin 与 Lite3 的真机部署问题，对统一闭环系统在真实平台上的实现与验证进行了系统阐述。